%==============================================================================
%  Homework template  --  999 Controls
%
%  To reuse for the next assignment:
%    1. copy this folder,  2. edit the METADATA block below,
%    3. edit the SOLUTIONS section,  4. run  `make`.
%
%  Build:  make          (regenerates code output, then builds main.pdf)
%          make clean    (removes build artifacts)
%==============================================================================
\documentclass[11pt,letterpaper]{article}

%------------------------------- METADATA -------------------------------------
\newcommand{\hwcourse}{EE 999 --- Introduction to Control Systems}
\newcommand{\hwcourseshort}{EE 999}  % running head, keep it short
\newcommand{\hwtitle}{Homework 2}
\newcommand{\hwauthor}{Patrick Gordon}
\newcommand{\hwdate}{\today}
\newcommand{\hwdue}{Due Wednesday, September 23, 2026}
%------------------------------------------------------------------------------

\usepackage[T1]{fontenc}
\usepackage[utf8]{inputenc}
\usepackage{textcomp}
\usepackage[letterpaper,margin=1in,headheight=14pt]{geometry}
\usepackage{amsmath,amssymb}
\usepackage{cancel}
\usepackage{array}
\usepackage{graphicx}
\usepackage{xcolor}
\usepackage{listings}
\usepackage{fancyhdr}
\usepackage{inconsolata}
\usepackage{pdflscape}
\usepackage[hidelinks]{hyperref}

%--- page furniture -----------------------------------------------------------
\pagestyle{fancy}
\fancyhf{}
\lhead{\hwauthor}
\chead{\hwtitle}
\rhead{\hwcourseshort}
\cfoot{\thepage}

%--- listing styles -----------------------------------------------------------
\definecolor{codebg}{HTML}{F7F7F4}
\definecolor{outbg}{HTML}{FCFCFC}
\definecolor{rule}{HTML}{C8C8C0}
\definecolor{kw}{HTML}{1C4FA1}
\definecolor{cmt}{HTML}{6A7A5A}
\definecolor{str}{HTML}{A0522D}

\lstdefinestyle{mcode}{
  language=Octave,
  basicstyle=\ttfamily\footnotesize,
  keywordstyle=\color{kw}\bfseries,
  commentstyle=\color{cmt}\itshape,
  stringstyle=\color{str},
  backgroundcolor=\color{codebg},
  frame=single,
  rulecolor=\color{rule},
  framesep=6pt,
  numbers=left,
  numberstyle=\ttfamily\scriptsize\color{gray},
  numbersep=8pt,
  showstringspaces=false,
  breaklines=true,
  columns=fullflexible,
  keepspaces=true,
  upquote=true,
}

\lstdefinestyle{moutput}{
  basicstyle=\ttfamily\footnotesize,
  backgroundcolor=\color{outbg},
  frame=single,
  rulecolor=\color{rule},
  framesep=6pt,
  showstringspaces=false,
  breaklines=true,
  columns=fullflexible,
  keepspaces=true,
  upquote=true,
}

%--- structure helpers --------------------------------------------------------
% \problem{1.1b}{Optional title}
\newcommand{\problem}[2]{%
  \bigskip
  \noindent{\large\bfseries Problem #1}%
  \ifx\relax#2\relax\else\ {\large --- #2}\fi
  \par\nobreak\vspace{2pt}\hrule\vspace{8pt}\nobreak
}

% \caplabel{Code}{hw_02_A.m}
\newcommand{\caplabel}[2]{%
  \par\nobreak\vspace{6pt}\nobreak
  \noindent{\small\bfseries #1}\ {\small\ttfamily \detokenize{#2}}%
  \par\nobreak\vspace{2pt}\nobreak
}

% \matlabcode{hw_02_A.m}  -- source listing pulled straight from the .m file
\newcommand{\matlabcode}[1]{%
  \caplabel{Code:}{#1}%
  \lstinputlisting[style=mcode]{#1}%
}

% \codeoutput{output/hw_02_A.txt} -- console output captured by the Makefile
\newcommand{\codeoutput}[1]{%
  \caplabel{Output:}{#1}%
  \IfFileExists{#1}{\lstinputlisting[style=moutput]{#1}}{\missing{#1}{run \texttt{make} to generate it}}%
}

% \answerfigure[width]{figures/foo.pdf}{Caption}  -- placeholder if not exported
% yet. Any format graphicx reads works: pdf, png, jpg.
% A script that plots is exported by the Makefile to figures/<script>.pdf, so
% \answerfigure{figures/hw_02_F.pdf}{...} is how a plotted answer is included.
% The optional argument is the printed width, default the full text width.
% Give a scan or screenshot a smaller one so it is not blown up past its own
% resolution:  \answerfigure[0.8\linewidth]{figures/some_scan.png}{...}
\newcommand{\answerfigure}[3][\linewidth]{%
  \begin{figure}[H]
    \centering
    \IfFileExists{#2}{\includegraphics[width=#1]{#2}}{\missing{#2}{export it, see README}}
    \caption{#3}
  \end{figure}%
}

%--- tables -------------------------------------------------------------------
% Centred paragraph column, so a description column wraps instead of running
% off the page:  \begin{tabular}{|c|P{0.5\linewidth}|}
\newcolumntype{P}[1]{>{\centering\arraybackslash}p{#1}}

% \begin{symboltable}{Signal}{Caption}  -- the symbol / units / description
% table used to define the variables of a problem. The first argument is the
% heading of the left column (Signal, Parameter, State, ...); leave the second
% empty for no caption.
%
%   \begin{symboltable}{Signal}{}
%     \sym{x(t)}{m}{Position of tire center}
%     \sym{r(t)}{m}{Road height (depends on $v_{car}(t)$)}
%   \end{symboltable}
%
% \sym puts its first argument in math mode and brackets the units, so write
% \sym{x(t)}{m}{...}, not \sym{$x(t)$}{[m]}{...}. Anything already math -- a
% subscript, a fraction -- goes in the description as usual: $\frac{1}{4}$.
%
% The optional argument is the width of the description column, for a table
% whose descriptions are much longer or much shorter than average:
%   \begin{symboltable}[0.35\linewidth]{Parameter}{}
\newcommand{\sym}[3]{$#1$ & [#2] & #3 \\ \hline}
\newenvironment{symboltable}[3][0.45\linewidth]{%
  \def\symtabcap{#3}%
  \begin{table}[H]
  \centering
  \renewcommand{\arraystretch}{1.4}
  \begin{tabular}{|c|c|P{#1}|}
  \hline
  #2 & Units & Description \\ \hline\hline
}{%
  \end{tabular}
  \ifx\symtabcap\empty\else\caption{\symtabcap}\fi
  \end{table}%
}

% \begin{answertable}{|c|c|c|}{Caption} -- any other table. Same placement and
% captioning as \answerfigure, but you write the rows yourself:
%
%   \begin{answertable}{|c|P{0.5\linewidth}|}{Pole locations.}
%     \hline
%     Pole & Behaviour \\ \hline\hline
%     $s=-2$ & decays in half a second \\ \hline
%   \end{answertable}
\newenvironment{answertable}[2]{%
  \def\anstabcap{#2}%
  \begin{table}[H]
  \centering
  \renewcommand{\arraystretch}{1.4}
  \begin{tabular}{#1}
}{%
  \end{tabular}
  \ifx\anstabcap\empty\else\caption{\anstabcap}\fi
  \end{table}%
}

% \begin{termtable}{4}{Caption} -- collect like powers of $s$ when multiplying
% out a characteristic polynomial (or any other expansion). The argument is the
% highest power; the columns and the $s^4 \dots s^0$ header are generated from
% it, so changing the degree is a one-character edit. Leave the second argument
% empty for no caption.
%
% Cells are already in math mode: write m_1m_2, not $m_1m_2$. One row per factor
% you multiplied out, one column per power; leave a cell empty where that factor
% contributes nothing to that power. Rows end with \\, like any tabular.
%
%   \begin{termtable}{4}{Terms of the characteristic polynomial.}
%     m_1m_2 & m_1b & m_1k_s       & bk_s       & k_s(k_s+k_w) \\
%            & m_2b & m_2(k_s+k_w) & b(k_s+k_w) &              \\
%            &      & -b^2         & -bk_s      & -k_s^2
%   \end{termtable}
%
% M is a centred column that is always in math mode. It uses \( \) rather than
% $ $ so that an empty cell stays empty -- a bare $$ would start display math.
\newcolumntype{M}{>{\(}c<{\)}}
\makeatletter
\newcount\tt@i
\newcommand{\tt@cols}{}
\newcommand{\tt@head}{}
\newcommand{\tt@build}[1]{%
  \def\tt@cols{}\def\tt@head{}%
  \tt@i=#1\relax
  \@whilenum\tt@i>\m@ne\do{%
    \edef\tt@cols{\tt@cols M}%
    \ifx\tt@head\@empty
      \edef\tt@head{s^{\the\tt@i}}%
    \else
      \edef\tt@head{\tt@head & s^{\the\tt@i}}%
    \fi
    \advance\tt@i\m@ne
  }%
}
\newenvironment{termtable}[2]{%
  \def\tt@cap{#2}%
  \tt@build{#1}%
  \begin{table}[H]%
  \centering
  \setlength{\tabcolsep}{12pt}%
  \renewcommand{\arraystretch}{1.3}%
  \edef\@tempa{\noexpand\begin{tabular}{\tt@cols}}\@tempa
  \tt@head \\ \hline
}{%
  \end{tabular}
  \ifx\tt@cap\@empty\else\caption{\tt@cap}\fi
  \end{table}%
}
\makeatother

% wideproblem environment -- a problem on its own landscape page, centred.
% Use for diagrams much wider than they are tall, where a portrait page would
% shrink the labels below readable size:
%   \begin{wideproblem}{1.1b}{Title}
%     \answerfigure{figures/x.pdf}{Caption}
%   \end{wideproblem}
\newenvironment{wideproblem}[2]{%
  \begin{landscape}%
    \problem{#1}{#2}%
    \vspace*{\fill}%
}{%
    \vspace*{\fill}%
  \end{landscape}%
}

% \prompt{What special property does b*d have?} -- restate the question
\newcommand{\prompt}[1]{%
  \noindent\textit{#1}\par\vspace{6pt}%
}

% \answer{...}  -- written answer. Leave the body empty for a visible
% placeholder box; fill it in and the box is replaced by the prose.
\newcommand{\answer}[1]{%
  \ifx\relax#1\relax\todo{write the answer here}\else#1\fi
  \par\vspace{4pt}%
}

\newcommand{\todo}[1]{%
  \fbox{\parbox{0.9\linewidth}{\centering\vspace{10pt}\rmfamily\itshape\small
    #1\vspace{10pt}}}%
}

\newcommand{\missing}[2]{%
  \fbox{\parbox{0.9\linewidth}{\centering\vspace{10pt}\ttfamily\small
    missing: \detokenize{#1}\\[4pt]\rmfamily\itshape #2\vspace{10pt}}}%
}

\usepackage{float}

%==============================================================================
\begin{document}

\begin{center}
  {\LARGE\bfseries \hwtitle}\\[4pt]
  {\large \hwcourse}\\[8pt]
  {\hwauthor} \quad\textbullet\quad {\hwdate}\\[4pt]
  {\small\itshape \hwdue}
\end{center}
\vspace{6pt}\hrule

%------------------------------------------------------------------------------
%  SOLUTIONS
%------------------------------------------------------------------------------

\clearpage
\problem{1.A.2.1}{What is a "free-body diagram"?}
\answer{A free body diagram is a diagram that shows the forces acting on an object or set of objects.}

\problem{1.A.2.2}{What are the two forms of Newton's law?}
\answer{Translational motion: $F=ma$; Rotational motion: $M = I \alpha$}

\problem{1.A.2.3}{For a structural process to be controlled, such as a robot arm, what is the meaning of "collocated control" and "Noncollocated control"?}
\answer{Collocated control means the sensor and actuator are rigidly attached to one another, whereas there is flexibility between the actuator and sensor in a noncollocated system.}

\problem{1.A.2.4}{State Kirchhoff's current law.}
\answer{Kirchoff's current law states that current entering any given electrical node must equal current exiting a node.}

\problem{1.A.2.5}{State Kirchhoff's voltage law.}
\answer{Kirchoff's voltage law states that the sum of all voltages around a closed circuit must equal zero.}

\clearpage
\problem{1.B}{Differential equation for mechanical system}
\answerfigure[0.8\linewidth]{figures/2_B_Schematic.pdf}{The two-mass system of Problem 1.B.}

\begin{symboltable}{Signal}{Signals of Problem 1.B.}
    \sym{x_1(t)}{m}{Position of mass $m_1$}
    \sym{x_2(t)}{m}{Position of mass $m_2$}
\end{symboltable}

\begin{symboltable}{Parameter}{Parameters of Problem 1.B.}
    \sym{m_1}{Kg}{mass of $m_1$}
    \sym{m_2}{Kg}{mass of $m_2$}
    \sym{K_{1A}}{N/m}{Outer spring stiffness, wall to $m_1$}
    \sym{K_{1B}}{N/m}{Outer spring stiffness, $m_2$ to wall}
    \sym{K_2}{N/m}{Inner spring stiffness}
    \sym{b_2}{N-s/m}{Shock absorber}
\end{symboltable}

\answer{Constituent equations: 
    \begin{align*}
        f_{K_{1A}}(t) &= K_{1A}x_1(t)\\
        f_{K_{1B}}(t) &= K_{1B}(-x_2(t))\\
        f_{K_2}(t) &= K_2(x_2(t)-x_1(t))\\
        f_{b_2} &= b_2(\dot{x}_2(t)-\dot{x}_1(t))\\
        f_{net_{m1}}(t) &= m_1\ddot{x}_1(t)\\
        f_{net_{m2}}(t) &= m_2\ddot{x}_2(t)
    \end{align*}
}

\answerfigure{figures/1_B_free_body.pdf}{Free body diagram for problem 1.B}

\answer{Continuity Equations:
    \begin{align*}
        K_1 &= K_{1A} = K_{1B}\\
        m &= m_1 = m_2\\
        \\
        m\ddot{x}_1(t) &= f_{K_2} + f_{b_2} - f_{K_{1A}} \\
                       &= K_2(x_2-x_1) + b_2(\dot{x}_2 - \dot{x}_1) - K_1x_1\\
                       &= K_2x_2 - K_2x_1 + b_2\dot{x}_2 - b_2\dot{x}_1 - K_1x_1\\
                       &= -b_2\dot{x}_1 - (K_1+K_2)x_1 + b_2\dot{x}_2 + K_2x_2 \\
        \\
        m\ddot{x}_2(t) &= f_{K_{1B}} - f_{K_2} - f_{b_2} \\
                       &= K_1( - x_2) - K_2(x_2 - x_1) - b_2(\dot{x}_2 - \dot{x}_1)\\
                       &= - K_1x_2 - K_2x_2 + K_2x_1 - b_2\dot{x}_2 + b_2\dot{x}_1\\
                       &= - b_2\dot{x}_2 - (K_1 + K_2)x_2 + b_2\dot{x}_1 + K_2x_1 \\
    \end{align*}
}

\answer{Standard form:
    \begin{align*}
        m\ddot{x}_1(t) + b_2\dot{x}_1(t) + (K_1 + K_2)x_1(t) &= b_2\dot{x}_2(t) + K_2x_2(t) + f(t)\\
        \\
        m\ddot{x}_2(t) + b_2\dot{x}_2(t) + (K_1 + K_2)x_2(t) &= b_2\dot{x}_1(t) + K_2x_1(t) \\
    \end{align*}

    *NOTE: $f(t)$ has been added to the $m_1$ equation so the transfer function could be created
}

\answer{Frequency domain:     
    \begin{align*}
        (ms^2 + b_2s + (K_1 + K_2))X_1(s) &= (b_2s + K_2)X_2(s) + F(s)\\
        \\
        (ms^2 + b_2s + (K_1 + K_2))X_2(s) &= (b_2s + K_2)X_1(s)
    \end{align*}

    Find $X_1(s)$ in terms of $X_2(s)$ so substitution may be performed:

    \[X_1(s) = \frac{ms^2 + b_2s + (K_1 + K_2)}{b_2s + K_2}X_2(s)\]

    Then perform the substitution:
    \begin{align*}
        (ms^2 + b_2s + (K_1 + K_2))\frac{ms^2 + b_2s + K_1 + K_2}{b_2s + K_2}X_2(s) &= (b_2s + K_2)X_2(s) + F(s) \\
        (ms^2 + b_2s + (K_1 + K_2))^2X_2(s) &= (b_2s + K_2)^2X_2(s) + (b_2s + K_2)F(s) \\
        ((ms^2 + b_2s + (K_1 + K_2))^2 - (b_2s + K_2)^2)X_2(s) &=  (b_2s + K_2)F(s) \\
    \end{align*}

    Next, gather the terms:

    \begin{termtable}{4}{}
       m^2 & mb_2 & m(K_1+K_2) &  &  \\
       & mb_2 & b_2^2 & b_2(K_1+K_2) &  \\
       &  & m(K_1+K_2) & b_2(K_1+K_2) & K_1^2 + 2K_1K_2 + K_2^2 \\
       &  & -b_2^2 & -b_2K_2 & \\
       &  &  & -b_2K_2 & -K_2^2
    \end{termtable}

    Then, combine them to write the entire transfer:

    \[ G_{FX_{2}(s)} = \frac{X_2(s)}{F(s)} = \frac{b_2s + K_2}{m^2s^4+ 2mb_2s^3 + 2m(K_1+K_2)s^2 + (2b_2(K_1 + K_2)-2b_2K_2)s + K_1^2 + 2K_1K_2} \]
}

\answer{Check Units:
    \begin{align*}
        \frac{\frac{Ns}{m}s^{-1} + \frac{N}{m}}{Kg^2s^{-4} + Kg\frac{Ns}{m}s^{-3} + Kg\frac{N}{m}s^{-2} + \frac{Ns}{m}\frac{N}{m}s^{-1} + (\frac{N}{m})^2} \\
        \\
        \frac{\frac{N\cancel{s}}{m}\cancel{s^{-1}} + \frac{N}{m}}{\frac{N^2\cancel{s^4}}{m^2}\cancel{s^{-4}} + \frac{N\cancel{s^2}}{m}\frac{N\cancel{s}}{m}\cancel{s^{-3}} + \frac{N\cancel{s^2}}{m}\frac{N}{m}\cancel{s^{-2}} + \frac{N\cancel{s}}{m}\frac{N}{m}\cancel{s^{-1}} + (\frac{N}{m})^2} \\
        \\
        \frac{\frac{N}{m}}{\frac{N^2}{m^2}} = \frac{m}{N}
    \end{align*}

    This makes sense because the input is a force and the output is a displacement, so the units of the transfer function should be distance over force.
}

\answer{The difference in velocities between $m_1$ and $m_2$ will decay because $b_2$ scales with velocity. Since there is no friction, the output $x_2(t)$ will not decay.}

\end{document}
